\section{PPO: Proximal Policy Optimization}

\subsection{Idea central: limitar el desplazamiento de la política}

Dos estrategias:

\textbf{Estrategia 1: restricción de KL}

$$\max_{\theta'} \tilde{J}(\theta') \quad \text{s.t.} \quad \mathbb{E}_{s \sim \pi_\theta}\left[D_{\mathrm{KL}}\left(\pi_{\theta'}(\cdot \mid s) \| \pi_\theta(\cdot \mid s)\right)\right] \leq \epsilon$$

Muy común en la optimización de preferencias de los LLM.

\textbf{Estrategia 2: Clipping}

En lugar de restringir directamente la política, se recortan los importance weights, eliminando el incentivo cuando la política se desvía en gran medida:

$$\tilde{J}(\theta') = \sum_{t,i} \min\left(\frac{\pi_{\theta'}(a_{i,t} \mid s_{i,t})}{\pi_\theta(a_{i,t} \mid s_{i,t})} \hat{A}^{\pi_\theta}, \; \mathrm{clip}\left(\frac{\pi_{\theta'}(a_{i,t} \mid s_{i,t})}{\pi_\theta(a_{i,t} \mid s_{i,t})}, 1-\epsilon, 1+\epsilon\right) \hat{A}^{\pi_\theta}\right)$$

\begin{importantbox}{El objetivo central de PPO}
Este es el surrogate objective final de PPO. Mediante el clipping y la operación de mínimo, la política ya no tiene incentivo para desviarse en gran medida de la política anterior, y se mantiene estable incluso al hacer actualizaciones de varios pasos.
\end{importantbox}

\subsection{Algoritmo completo de PPO}

\begin{enumerate}
    \item Recolectar un lote de datos con la política actual $\pi_\theta$
    \item Ajustar la función de valor $\hat{V}^\pi_\phi$
    \item Calcular la ventaja $\hat{A}^{\pi_\theta}$
    \item Sobre este lote de datos, hacer varios pasos de actualización de gradiente del objetivo de PPO (aproximadamente 3--10 pasos)
    \item Repetir
\end{enumerate}

\subsection{Resumen de la sección}

PPO limita el desplazamiento de la política recortando los importance weights, lo que permite hacer varias actualizaciones sobre el mismo lote de datos, y es uno de los algoritmos de RL on-policy más utilizados en la actualidad.

